\chapter{Past Paper: Sample Examination}

\section{Paper metadata}

\begin{tabularx}{\textwidth}{@{}>{\bfseries\sffamily\color{CourseNavy}}L{38mm}Y@{}}
Academic year & 2025--2026 \\
Assessment & Final examination \\
Duration & 120 minutes \\
Total marks & 100 \\
Permitted materials & Non-programmable calculator; one formula sheet \\
Source status & Demonstration question written for this template \\
Official solution & Not applicable \\
\end{tabularx}

\begin{editorialnote}[Source labelling]
Do not remove the year, assessment type, page number, question number, or source status
from archived questions. If a question is reconstructed from memory or an incomplete
scan, label it explicitly as a reconstruction.
\end{editorialnote}

\section{Question 1}

\begin{pastpaperproblem}{Linear systems and interpretation}{exam-q1}
\probleminfo{Linear systems}{Standard}{Sample Examination, Q1}{12}
Consider
\[
  A=\begin{pmatrix}1&2\\2&4\end{pmatrix},
  \qquad
  \bm{b}=\begin{pmatrix}3\\6\end{pmatrix}.
\]
\begin{enumerate}
  \item Determine the rank of $A$.
  \item Solve $A\bm{x}=\bm{b}$.
  \item Explain why the solution is not unique.
\end{enumerate}
\end{pastpaperproblem}

\begin{markscheme}
\begin{tabularx}{\textwidth}{@{}L{18mm}Y C{18mm}@{}}
\toprule
Part & Expected evidence & Marks \\
\midrule
(a) & Correct row reduction or identification of dependent rows; conclusion $\rank A=1$. & 3 \\
(b) & One free parameter and a correct parametrisation of all solutions. & 6 \\
(c) & Correct reference to a nontrivial null space or dependent columns. & 3 \\
\bottomrule
\end{tabularx}
\end{markscheme}

\begin{coursesolution}[Complete solution]
The second row of $A$ is twice the first, so $\rank A=1$.
The system reduces to the single equation
\[
  x+2y=3.
\]
Let $y=t$. Then $x=3-2t$, and therefore
\[
  \bm{x}=\begin{pmatrix}3-2t\\t\end{pmatrix},
  \qquad t\in\R.
\]
The solution is not unique because the null space contains the nonzero vector
$(-2,1)^{\mathsf T}$. Equivalently, the coefficient matrix does not have full column
rank.
\end{coursesolution}

\section{Cross-paper classification}

When several papers have been entered, maintain a summary table such as the following.

\begin{tabularx}{\textwidth}{@{}Y C{20mm} C{22mm} C{24mm}@{}}
\toprule
Topic & Frequency & Typical marks & Main format \\
\midrule
Linear systems & High & 8--15 & Calculation + interpretation \\
Eigenvalues & Medium & 10--18 & Calculation + short proof \\
Vector spaces & High & 10--20 & Definitions + proof \\
\bottomrule
\end{tabularx}
